\subsection{Linear compactness properties of complete filtered modules}

Recall that, if E is an A-module, an affine subset (or an affine linear variety) of E is any subset F which is empty or of the form \(a + M\) , where \(a \in E\) and M is a submodule of E called the direction of F (Algebra, Chapter II, § 9, nos. 1 and 3).

PROPOSITION 7. Let A be a filtered ring, E a filtered A-module and \((\mathbf{E}_{n})\) the filtration of E; suppose that \(E_{0} = E\) , that the \(E_{n}\) are submodules of E, that the A-modules \(E/E_{n}\) are Artinian and finally that the topological group E is Hausdorff and complete. Then the intersection of a decreasing sequence of non-empty closed affine subsets of E is non-empty.

We have seen in no. 6 that, since E is Hausdorff and complete, it is identified with \(\tilde{\mathbf{E}} = \varprojlim \mathrm{E} / \mathrm{E}_n\) . Let \((\mathbf{W}_p)\) be a decreasing sequence of non-empty closed affine subsets of E and, for all \(n \geqslant 0\) , let \(\mathbf{W}_{p,n}\) be the canonical image of \(\mathbf{W}_p\) in \(\mathrm{E} / \mathrm{E}_n\) ; we are going to construct a sequence \(x = (x_n) \in \tilde{\mathbf{E}}\) such that \(x_n \in \mathbf{W}_{p,n}\) for all \(p\) and all \(n\) ; hence \(x \in \mathbf{W}_p + \mathbf{E}_n\) for all \(p\) and all \(n\) and, as the \(\mathbf{W}_p\) are closed, \(x \in \mathbf{W}_p\) for all \(p\) (no. 5), which will prove the proposition.

As \(E/E_{0}=0\) , we shall take \(x_{0}=0\) . Suppose that the \(x_{i}\) are defined for \(0\leqslant i\leqslant n-1\) and let \(W_{p,n}^{\prime}\) be the set of elements of \(W_{p,n}\) whose canonical image in \(E/E_{n-1}\) is \(x_{n-1}\) ; as \(x_{n-1}\in W_{p,n-1}\) and \(W_{p,n-1}\) is the canonical image of \(W_{p,n}\) , \(W_{p,n}^{\prime}\) is non-empty and is obviously an affine subset of \(E/E_{n}\) ; moreover the sequence \((W_{p,n}^{\prime})_{p\in\mathbf{N}}\) is decreasing. As \(E/E_{n}\) is Artinian, this sequence is stationary (otherwise the sequence of submodules of \(E/E_{n}\) which are the directions of the \(W_{p,n}^{\prime}\) would be strictly decreasing, which is absurd). It is sufficient then to take \(x_{n}\) to be an element of \(\bigcap_{p\in\mathbf{N}}W_{p,n}^{\prime}\) and the construction of \((x_{n})\) can then be performed by induction.

PROPOSITION 8. Suppose that A and E satisfy the hypotheses of Proposition 7. Let \((\mathbf{F}_{p})\) be a decreasing sequence of closed submodules of E such that \(\bigcap_{p} F_{p} = 0\) . Then, for every neighbourhood V of 0 in E, there exists p such that \(F_{p} \subset V\) (in other words, the base of the filter \((\mathbf{F}_{p})\) converges to 0).

We may assume that V is one of the \(E_{n}\) , in which case E/V is Artinian. Let us write \(F_{p}^{\prime} = (F_{p} + V)/V\) ; as the \(F_{p}^{\prime}\) form a decreasing sequence of submodules of E/V, there exists an integer j such that \(F_{p}^{\prime} = F_{j}^{\prime}\) for all \(p \geqslant j\) . We shall see that \(F_{j}^{\prime} = \{0\}\) , which will complete the proof. Let \(x \in F_{j}^{\prime}\) and let \(W_{p}\) be the set of elements of \(F_{p}\) whose image in E/V is \(x (p \geqslant j)\) ; by definition of j, the \(W_{p}\) are non-empty closed affine subsets of E and obviously \(W_{p+1} \subset W_{p}\) ;

then it follows from Proposition 7 that there exists an element \(y\) belonging to all the \(W_p\) . As \(W_p \subset F_p\) and \(\bigcap_{p \in \mathbf{N}} F_p = \{0\}\) , \(y = 0\) ; since \(x\) is the canonical image of \(y\) in \(E / V\) , \(x = 0\) (cf.~Exercises 15 to 21).

